% uid: ""
% school: "lycee"
% level: "3eme"
% section: ["Economie et Gestion"]
% chapter: "Suites réelles"
% title: "Placement à intérêts simples"
% difficulty: "Facile"
% points: null
% tags: ["suites", "suite-arithmetique", "interets-simples", "terme-general", "seuil"]
% provenance: "own"
% dressing: "suite arithmétique"
% rung: 1
% batch: "2026-11-20-suites-3eco-A"

% ---STATEMENT---
Une personne place $2\,000$ dinars à intérêts simples au taux annuel de $4\,\%$ : chaque année, l'intérêt est calculé sur le capital de départ. On note $U_n$ le capital (en dinars) au bout de $n$ années, donc $U_0=2\,000$.

\begin{enumerate}
  \item Calculer $U_1$ et $U_2$. Montrer que $(U_n)$ est une suite arithmétique dont on donnera la raison.
  \item Exprimer $U_n$ en fonction de $n$. Calculer $U_{10}$.
  \item À partir de quelle année le capital dépasse-t-il $3\,000$ dinars ?
\end{enumerate}

% ---SOLUTION---
\begin{enumerate}
  \item L'intérêt annuel est toujours calculé sur $2\,000$ : $4\,\%\times2\,000=80$ dinars.
  $$U_1=2\,000+80=2\,080\qquad U_2=2\,080+80=2\,160.$$
  Pour tout $n$, $U_{n+1}=U_n+80$ : on ajoute toujours le même nombre. $(U_n)$ est donc arithmétique de raison $\ib{$r=80$}$.
  \item \textbf{Rappel.}\par
  \begin{prop}[Terme général d'une suite arithmétique]
  $u_n=u_0+n\,r$.
  \end{prop}
  $$ \ib{$U_n=2\,000+80\,n$}\qquad U_{10}=2\,000+800=2\,800\text{ dinars}.$$
  \item On résout l'inéquation $U_n>3\,000$ :
  $$2\,000+80\,n>3\,000\iff80\,n>1\,000\iff n>12{,}5.$$
  Le plus petit entier est $n=13$. \emph{Vérification} : $U_{12}=2\,960<3\,000$ et $U_{13}=3\,040>3\,000$ $\checkmark$

  Le capital dépasse $3\,000$ dinars à partir de la $13^{\text{e}}$ année.

  \textbf{Piège :} à intérêts \emph{composés} le capital serait multiplié par $1{,}04$ chaque année : suite géométrique, croissance plus rapide.
\end{enumerate}

\textbf{Idée :} montant constant chaque année : on ajoute toujours le même nombre, c'est une suite arithmétique.
